%% CloudKC-Bench — IEEE journal draft (IEEEtran).
%% Target: IEEE TIFS / Computers & Security. Compiles on Overleaf as-is.
%% Ownership (per plan v3 §11): Atishay = Abstract, System, Results, Failure Analysis,
%% Limitations. Anna = Introduction, Related Work, Discussion, Conclusion.
%% [RESULT] = fill from `benchmark.cli analyze/failures` once the run completes.
\documentclass[journal]{IEEEtran}

\usepackage{cite}
\usepackage{amsmath,amssymb}
\usepackage{graphicx}
\usepackage{booktabs}
\usepackage{xcolor}
\usepackage{array}
\usepackage{url}
\usepackage[hidelinks]{hyperref}

% Placeholders — make unfilled results impossible to miss.
\newcommand{\result}[1]{\textcolor{red}{\textbf{[RESULT: #1]}}}
\newcommand{\todo}[1]{\textcolor{blue}{[TODO: #1]}}
\newcommand{\cmark}{\ensuremath{\checkmark}}
\newcommand{\xmark}{\ensuremath{\times}}

% Self-healing figure: drop in the generated PDF if it exists, else a labelled
% placeholder box so the draft always compiles. The moment the PDF lands in
% paper/figures/ the real figure appears with no LaTeX edits. Data figures are
% produced by `python -m benchmark.plots` (see benchmark/plots.py).
\newcommand{\autofig}[2]{%
  \IfFileExists{#1}{\includegraphics[width=\linewidth]{#1}}{%
    \setlength{\fboxsep}{10pt}%
    \fbox{\parbox[c][3cm][c]{0.88\linewidth}{\centering
      \textcolor{red}{\textbf{[FIGURE PLACEHOLDER]}}\\[3pt]
      \texttt{\small\detokenize{#1}}\\[5pt]  % detokenize: underscores in filenames stay literal
      \footnotesize\itshape #2}}}%
}

\begin{document}

\title{CloudKC-Bench: An Open Benchmark and Controlled Study of LLM
Cross-Domain Correlation for AWS Multi-Stage Kill-Chain Reconstruction}

\author{Atishay~Jain,~Anna~M.,~and~S.~Nagasundari%
\thanks{The authors are with the Centre for Information Security, Forensics and
Cyber Resilience (ISFCR), PES University, Bengaluru, India.}%
\thanks{Corresponding author: Atishay Jain (e-mail: atishay.offical@gmail.com).}}

\markboth{Preprint --- draft \today}{Jain \MakeLowercase{\textit{et al.}}: CloudKC-Bench}

\maketitle

\begin{abstract}
Multi-stage attacks on cloud infrastructure leave traces scattered across
heterogeneous control-plane and network telemetry---IAM/CloudTrail API events, VPC
Flow Logs, S3 access logs, and EC2 lifecycle events---such that no single source
contains a whole attack. Whether Large Language Models (LLMs) can correlate across
these sources to reconstruct the kill chain, and whether they do so better than
simpler baselines, is an open and largely unmeasured question. We present
\emph{CloudKC-Bench}, the first open benchmark targeting AWS control-plane
multi-stage kill-chain reconstruction: a taxonomy-structured set of
$\sim$70 scenarios grounded in documented real incidents and ATT\&CK-for-Cloud
techniques, each released with a machine-checkable ground-truth manifest and a
published mechanical matching function. Using the benchmark we run a pre-registered
four-arm ablation (multi-agent LLM, single-agent LLM, single agent without a
pre-filter, and a deterministic rules-only correlator) plus an independent community
Sigma-rules baseline, holding model, telemetry, and compute fixed across arms.
We report standard detection metrics per scenario category with bootstrap confidence
intervals and effect sizes, and resolve two pre-registered hypotheses. On a real-AWS run
(qwen2.5:32b), the LLM performs cross-domain attack-\emph{event} detection at higher
precision than the rules (up to $1.00$ vs $0.65$ on multi-stage kill chains), but is
markedly worse at exact ATT\&CK-\emph{technique} attribution (recall $\sim$0.12 vs
$0.87$); multi-agent decomposition adds no measurable value over a single agent.
We further quantify the cost/latency of the deterministic pre-filter and give a
failure-mode taxonomy showing when and why reconstruction fails. All artifacts---the
generator, manifests, schema, matching function, and analysis scripts---are released
to enable independent replication.
\end{abstract}

\begin{IEEEkeywords}
Cloud security, threat detection, large language models, benchmark, MITRE ATT\&CK,
kill-chain reconstruction, AWS, multi-agent systems.
\end{IEEEkeywords}

% =====================================================================
\section{Introduction}\label{sec:intro}
% Owner: Anna
Cloud control-plane telemetry is siloed by design: a single attack stage such as
privilege escalation surfaces in CloudTrail, while the ensuing exfiltration surfaces
in S3 access logs \emph{and} VPC Flow Logs. Reconstructing a multi-stage kill chain
therefore requires correlation \emph{across} telemetry sources that share no join key
and arrive with different, variable delivery latencies. LLMs have been proposed for
security triage, but claims about their value are frequently made without fair
baselines, without reproducible data, and without honest treatment of
non-determinism---exactly the flaws that sink security-ML papers.

\textbf{Problem.} Given heterogeneous AWS control-plane and network telemetry produced
during a multi-stage attack, reconstruct the kill chain---identify which events belong
to the attack, in what order, mapped to which ATT\&CK-for-Cloud techniques---and do so
measurably better, or measurably no better, than fair baselines.

\textbf{Gap.} There is no public AWS control-plane multi-stage kill-chain dataset with
machine-checkable ground-truth manifests (Section~\ref{sec:related}). Existing cloud
security benchmarks target a different platform (Azure), a different task (malware
analysis, threat-intel reasoning, architectural threat modeling, insider threat), or
produce detection rules rather than reconstructed chains.

\textbf{Contributions.}
\begin{itemize}
\item \textbf{C1.} \emph{CloudKC-Bench}, the first open benchmark targeting AWS
control-plane multi-stage kill chains, with ground-truth manifests and a published
mechanical matching function (Section~\ref{sec:bench}).
\item \textbf{C2.} A pre-registered four-arm ablation on identical telemetry that
attributes any detection advantage to a specific layer---pre-filter, LLM, or
multi-agent decomposition (Sections~\ref{sec:setup},~\ref{sec:results}).
\item \textbf{C3.} A first-class cost/latency analysis of deterministic pre-filtering,
with measured per-category filtering ratios and a cost--accuracy trade-off.
\item \textbf{C4.} A failure-mode taxonomy: which categories, techniques, and
false-positive patterns break reconstruction (Section~\ref{sec:failure}).
\item \textbf{C5.} A reproducibility artifact: pinned model, frozen scenarios, released
code, manifests, and analysis scripts.
\end{itemize}
We deliberately state H2 as a null and report negative results honestly; the benchmark
stands as a contribution regardless of whether the system ``wins''.

% =====================================================================
\section{Background}\label{sec:background}
% Owner: Atishay
\subsection{AWS control-plane telemetry}
CloudSentinel/CloudKC-Bench considers four sources. \emph{CloudTrail} records every
API call (\texttt{eventName}, \texttt{userIdentity}, \texttt{sourceIPAddress},
\texttt{eventTime}, request/response elements); its S3 delivery latency is typically
5--15 minutes. \emph{VPC Flow Logs} record accepted/rejected IP traffic at the ENI
level (\texttt{srcaddr}, \texttt{dstaddr}, ports, protocol, action, bytes).
\emph{S3 access logs / data events} record object operations (\texttt{GetObject},
\texttt{PutObject}, \texttt{DeleteObject}, \texttt{CopyObject}). \emph{EC2 lifecycle}
events record launches, terminations, and image/snapshot operations.

\subsection{Kill chains and ATT\&CK for Cloud}
A multi-stage cloud attack follows a recognisable progression (initial access,
persistence, privilege escalation, defense evasion, credential access, discovery,
lateral movement, exfiltration, impact), each stage typically visible in a different
telemetry source. We ground every scenario stage in an ATT\&CK-for-Cloud technique
identifier so that realism is externally citable rather than asserted.

\subsection{LLMs for security triage}
Prior work applies LLMs to explanation, rule generation, and QA over security logs.
Our position is that the multi-agent architecture is an \emph{instrument}, not a
contribution: its value is tested by ablation, and a null result is a valid finding.

% =====================================================================
\section{Related Work}\label{sec:related}
% Owner: Anna. Coverage-gap assessment from docs/week1/09 (papers read, Jul 2026).
We organise prior work as a taxonomy and, for each theme, state what it leaves open.

\subsection{LLM-assisted security detection}
LLMCloudHunter~\cite{schwartz2025llmcloudhunter} generates Sigma detection-rule
candidates from cyber-threat-intelligence reports for cloud, and HuntGPT~\cite{ali2023huntgpt}
couples a random-forest network IDS on KDD99 with GPT-3.5 for explanation.
\emph{Leaves open:} both target rule generation or single-host anomaly explanation, not
reconstruction of multi-stage AWS control-plane kill chains with mechanical scoring.

\subsection{Cloud-native threat detection}
Managed services (e.g., GuardDuty) and attack-emulation tooling such as Stratus Red
Team~\cite{stratus} and CloudGoat~\cite{cloudgoat} define ATT\&CK-mapped AWS techniques.
\emph{Leaves open:} these are detectors or attack generators, not open benchmarks with
ground-truth manifests and a published matching function.

\subsection{Multi-agent architectures for security}
Agentic and provenance-based hunting~\cite{mukherjee2025provenance} report gains from
RAG and agent orchestration on host-provenance graphs.
\emph{Leaves open:} host endpoint forensics, not cloud control-plane correlation; the
value of decomposition is asserted, rarely ablated---our H2 tests it directly.

\subsection{Security benchmarks: design, coverage, and gaps}
Table~\ref{tab:coverage} positions CloudKC-Bench against the five closest benchmarks.
CyberSOCEval~\cite{deason2025cybersoceval} evaluates malware analysis and threat-intel
reasoning; ExCyTIn-Bench~\cite{wu2025excytin} builds an \emph{Azure} threat-investigation
benchmark scored as QA over investigation graphs; OrgForge-IT~\cite{orgforge2026} is a
verifiable synthetic \emph{insider-threat} benchmark; and ACSE-Eval~\cite{acseeval2025}
evaluates LLM threat \emph{modeling} of AWS architecture/IaC. No competitor satisfies all
six properties simultaneously.
\emph{Leaves open:} an AWS-control-plane, telemetry-driven, machine-checkable
kill-chain-reconstruction benchmark---the gap CloudKC-Bench fills.

\begin{table*}[t]
\caption{Coverage gap. \cmark~yes, \xmark~no, $\circ$~partial. Assessed from the
papers (Jul 2026); partial cells to be confirmed against full text.}
\label{tab:coverage}
\centering
\begin{tabular}{lcccccc}
\toprule
Property & CyberSOCEval & ExCyTIn & OrgForge-IT & LLMCloudHunter & ACSE-Eval & \textbf{Ours} \\
\midrule
AWS control-plane specific      & \xmark & \xmark & \xmark & $\circ$ & $\circ$ & \cmark \\
Multi-stage kill chains         & \xmark & \cmark & \xmark & \xmark  & \xmark  & \cmark \\
Open + reproducible             & \cmark & \cmark & \cmark & \xmark  & \cmark  & \cmark \\
Machine-checkable ground truth  & $\circ$ & \cmark & \cmark & \xmark & $\circ$ & \cmark \\
ATT\&CK-for-Cloud mapped        & \xmark & \xmark & \xmark & $\circ$ & $\circ$ & \cmark \\
Real-AWS validated              & \xmark & \xmark & \xmark & \xmark  & \xmark  & \cmark \\
\bottomrule
\end{tabular}
\end{table*}

\subsection{Kill-chain and ATT\&CK modeling}
MITRE ATT\&CK for Cloud~\cite{mitreattack} provides the technique taxonomy we use to
label stages. \emph{Leaves open:} a scored reconstruction task over live telemetry.

% =====================================================================
\section{CloudKC-Bench Design}\label{sec:bench}
% Owner: Atishay
\subsection{Scenario taxonomy}
The benchmark comprises $\sim$70 scenarios across five never-pooled categories:
single-domain attacks (12; a sanity floor), multi-stage kill chains (15; the primary
test), low-and-slow evasion (12), ephemeral compute (10), and benign noise (10;
false-positive measurement), plus a sealed held-out set (10). Every scenario cites a
documented real incident pattern (e.g., the Capital One 2019 SSRF$\rightarrow$IMDS
$\rightarrow$S3 chain~\cite{capitalone2019}, SCARLETEEL~\cite{scarleteel2023}, TeamTNT
cryptojacking~\cite{teamtnt2020}) and one or more ATT\&CK-for-Cloud techniques.

\subsection{Ground-truth manifests and the answer-key boundary}
Each scenario ships a machine-checkable manifest: an ordered list of stages, each with a
technique id, telemetry source, evidence event ids, and a timestamp range, validated
against a published JSON Schema. Crucially, ground truth lives \emph{only} in the
manifest and an \texttt{is\_ground\_truth} flag---never in the event payload the systems
read---so an attack API call is, in isolation, indistinguishable from benign activity.

\subsection{Mechanical matching function}
Scoring is an exact, deterministic function of (reconstructed chain, manifest); no human
or LLM judging. A reported stage matches a ground-truth stage iff the technique id and
telemetry source match \emph{and} the cited evidence overlaps the ground-truth evidence
(preventing ``right answer, wrong reason'' credit). Matching is one-to-one; partial credit
is awarded per stage; order is scored via a documented penalty. Benign scenarios have no
stages, so any reported stage is a false positive.

\subsection{Quality controls}
Real-incident grounding, inter-rater review of $\geq$20\% of scenarios, and an authorship
separation for the held-out set (sealed before system finalisation) address the
synthetic-realism and selection-bias threats. A measured cross-service clock model
documents delivery lag/skew so temporal claims are defensible.

% =====================================================================
\section{System Architecture}\label{sec:system}
% Owner: Atishay. The architecture is the instrument, not the contribution.
Figure~\ref{fig:pipeline} shows the reference pipeline (arm A1) and where the ablation
arms diverge from it.
Telemetry is normalised into a SQLite state cache. A deterministic, instrumented
pre-filter reduces event volume before any LLM call, recording events-in/out and
pre-filter recall. The full system (arm A1) distributes filtered events to four domain
hunters (CloudTrail, VPC, S3, EC2), whose candidate detections a coordinator correlates
into an ordered chain. All arms emit the identical reconstructed-chain schema, so the
comparison is not confounded by output format. The LLM backend is pluggable---a
deterministic offline backend for reproducible pipeline tests, a local Ollama model, or a
hosted API (Gemini/OpenAI-compatible)---and is pinned per experiment.

\begin{figure}[t]
\centering
% arms_pipeline.svg exists in paper/figures/; export it to PDF for pdflatex
% (Overleaf: right-click the .svg -> "compile as PDF", or use inkscape/rsvg-convert).
\autofig{figures/arms_pipeline.pdf}{Export paper/figures/arms\_pipeline.svg to PDF.}
\caption{CloudKC-Bench arm pipeline. Normalised telemetry $\rightarrow$ deterministic
pre-filter $\rightarrow$ (A1) four domain hunters + coordinator. A2 replaces the hunters
with a single agent; A3 removes the pre-filter; A4 replaces the LLM with deterministic
rules. All arms emit the identical reconstructed-chain schema.}
\label{fig:pipeline}
\end{figure}

% =====================================================================
\section{Experimental Setup}\label{sec:setup}
% Owner: Atishay
\subsection{Arms and baselines}
Four arms hold model, telemetry, and compute fixed and vary one factor each:
\textbf{A1} pre-filter + multi-agent; \textbf{A2} pre-filter + single agent
(A1 vs A2 isolates decomposition, H2); \textbf{A3} single agent on raw logs
(A2 vs A3 isolates the pre-filter, C3); \textbf{A4} pre-filter + deterministic rules,
no LLM (A1/A2 vs A4 isolates the LLM, H1). An independent community Sigma-rules
baseline (\textbf{SIGMA}) guards against the objection that A4 was designed to lose.

LLMCloudHunter~\cite{schwartz2025llmcloudhunter} generates detection \emph{rules} from
CTI reports rather than reconstructing kill chains from live telemetry; comparing it on
CloudKC-Bench would be a task mismatch, so it is positioned in the related-work taxonomy
(Section~\ref{sec:related}) rather than run as a benchmark arm. We did enable GuardDuty on
the real-AWS account and collect its findings, but it produced no detections mapping to the
emulated techniques (Section~\ref{sec:threats}); we therefore exclude it as a \emph{scored}
baseline and retain SIGMA as the independent rules comparison, which needs no managed service.

\subsection{Environments}
Scenarios run in two environments emitting the identical schema: a reproducible layer
(deterministic synthetic generation / LocalStack), and---budget-gated---a real-AWS
sandbox for external validity. Real-AWS results, where available, are the headline;
the reproducible layer is the consistency check.

\subsection{Pre-registration and power}
The research question, H1's required margin ($\geq 0.15$ absolute multi-stage recall),
H2's equivalence band, and all tests were fixed in writing before any result. An a-priori
power analysis (paired design, $\alpha=0.05$, power $0.80$) shows that after
Holm--Bonferroni over the full comparison grid no category reaches the corrected $n$ for
moderate effects; we therefore pre-register \emph{two} primary confirmatory tests
(H1 and H2 on the multi-stage category) and treat all other comparisons as exploratory,
reported with effect sizes and bootstrap confidence intervals. The reported run pins
\texttt{qwen2.5:32b-instruct-q4\_k\_m} (local Ollama) over the dev set, 3 seeds, on real-AWS telemetry.

\subsection{Metrics and protocol}
Primary metrics are TPR/recall, FPR, precision, F1, and AUC where meaningful, reported
\emph{per category}. The unit of analysis is the scenario; seeds are repeated measures
aggregated to the scenario level. Effect sizes (Cohen's $d_z$) and bootstrap 95\%
confidence intervals are primary; $p$-values (paired permutation test) are secondary and
Holm--Bonferroni corrected.

\textbf{TCP disjoint-stream robustness test (DoD item 11).}
The correlator assigns stage order by sorting candidates on \texttt{event\_time};
different telemetry sources carry source-specific delivery lags (CloudTrail $\sim$7 min,
VPC $\sim$75 s, EC2 $\sim$30 s). We sweep a per-source timestamp offset $\Delta T \in
[-900, +900]$ s and measure the minimum $|\Delta T|$ at which \texttt{order\_penalty}>0
(\emph{robustness threshold}). This threshold is reported alongside the clock-model skews
(Section~\ref{sec:background}) to justify the correlation window choice. Crucially,
recall is invariant to $\Delta T$ (the matching function does not use timestamps), so a
rising order penalty represents a pure ordering failure rather than a detection failure.
Run \texttt{python -m benchmark.cli tcp-robustness} to reproduce.

% =====================================================================
\section{Results}\label{sec:results}
% Owner: Atishay.
\noindent\textbf{Run configuration.} The numbers below are from a \textbf{real-AWS} run of the full dev set
(59 scenarios across five categories, 3 seeds) with \textbf{qwen2.5:32b} (served locally via Ollama) as the LLM;
telemetry was captured by executing each scenario against a real AWS sandbox account. The held-out set remains
sealed. \todo{add the qwen2.5:7b/synthetic run as a model-scale + reproducibility comparison row.}

\subsection{Two scoring lenses}
We report reconstruction under two lenses (Section~\ref{sec:bench}): \emph{technique-level} (a stage is
correct only if the ATT\&CK technique, telemetry source, and cited evidence all match) and \emph{event-level}
(technique-agnostic: did the arm cite the ground-truth attack \emph{events} at all?). Table~\ref{tab:twolens}
and Fig.~\ref{fig:twolens} show they tell very different stories.

\begin{table}[t]
\caption{Multi-stage kill chains: two scoring lenses (qwen2.5:32b, real-AWS, $n{=}15$, 3 seeds).
Technique-level uses parent-technique credit; event-level is technique-agnostic.}
\label{tab:twolens}
\centering
\begin{tabular}{lcccc}
\toprule
 & A1 & A2 & A3 & \textbf{A4 (rules)} \\
\midrule
Technique-level recall      & 0.117 & 0.078 & 0.146 & \textbf{0.867} \\
Event-level recall          & 0.894 & 0.762 & 0.816 & 0.967 \\
Event-level precision       & 0.815 & 0.952 & \textbf{1.000} & 0.649 \\
Tokens / scenario           & 1814  & 1314  & 1393  & 0 \\
\bottomrule
\end{tabular}
\end{table}

\noindent The LLM arms \textbf{detect attack events precisely} --- event-level recall $0.76$--$0.89$ at
precision $0.82$--$1.00$, \emph{higher precision than the rules} ($0.65$; A3 reaches $1.00$) --- yet score
near-zero at the technique level. The rules (A4) win technique-level recall ($0.867$) because they encode the
exact sub-technique each scenario uses; the LLM, even given the closed technique vocabulary and parent-technique
credit, rarely names the exact technique. The bottleneck is \textbf{technique attribution, not cross-domain
event detection}.

\begin{figure}[t]
\centering
\autofig{figures/two_lens.pdf}{Generate: \texttt{benchmark.plots --results <run>.csv
--detection <det>.csv --category multi\_stage\_kill\_chain}.}
\caption{The two-lens gap on multi-stage kill chains: per arm, technique-level recall
(strict) vs event-level recall (technique-agnostic). The LLM arms (A1--A3) find the
attack events but rarely name the exact technique; A4 rules do both because their
signatures encode the answer key (real-AWS, qwen2.5:32b, $n{=}15$, 3 seeds).}
\label{fig:twolens}
\end{figure}

\begin{figure}[t]
\centering
\autofig{figures/recall_by_category.pdf}{Generate: \texttt{benchmark.plots --results <run>.csv}.}
\caption{Technique-level recall by arm across the five categories (qwen2.5:32b,
real-AWS). A4 rules dominate the strict metric everywhere; the LLM arms are low
across the board, consistent with the attribution bottleneck.}
\label{fig:recallcat}
\end{figure}

\subsection{Primary hypotheses}
\textbf{H1 (LLM vs rules, technique-level recall):} the best LLM arm (A1) trails A4 by a mean of $-0.750$
(95\% CI $[-0.841,-0.654]$, Cohen's $d_z=-3.91$, permutation $p<0.001$, Holm-corrected). H1 is
\textbf{not supported} --- and, on the strict technique-level metric, the LLM is significantly \emph{worse}.
The same ranking holds across all five categories (Fig.~\ref{fig:recallcat}).
\textbf{H2 (multi-agent vs single, F1):} A1 vs A2 differ by $+0.027$ (95\% CI $[-0.051,0.118]$, $d_z=0.16$,
$p=0.59$): we \textbf{fail to reject} H2 --- multi-agent decomposition adds no measurable value, while costing
$\sim$38\% more tokens (1814 vs 1314).

\subsection{Pre-filter cost and latency (C3)}
The deterministic pre-filter keeps $\sim$90\% of events on multi-stage (filtering ratio $0.90$; A3 runs
unfiltered at $1.0$). It reduces LLM token cost (A2 $1314$ vs A3 $1393$, $\approx$6\%) --- but at an
event-recall cost: \emph{A3 (no pre-filter) detects more attack events than A2} ($0.816$ vs $0.762$). Thus the
pre-filter is a genuine cost-vs-recall trade-off, and for a capable reader model it discards real detections.
A4 (rules) and SIGMA incur zero LLM cost (Fig.~\ref{fig:costacc}).

\begin{figure}[t]
\centering
\autofig{figures/cost_accuracy.pdf}{Generate: \texttt{benchmark.plots --results <run>.csv
--category multi\_stage\_kill\_chain}.}
\caption{Cost--accuracy trade-off on multi-stage kill chains: tokens per scenario vs
technique-level recall, per arm (qwen2.5:32b, real-AWS). A4/SIGMA sit at zero token
cost; the LLM arms pay tokens without a technique-level return on this metric.}
\label{fig:costacc}
\end{figure}

% =====================================================================
\section{Failure-Mode Analysis}\label{sec:failure}
% Owner: Atishay.
\textbf{Technique attribution is the dominant LLM failure.} Because event-level precision is high
($0.82$--$1.00$ on multi-stage) while technique-level recall is only $\sim$0.12, almost every event the LLM
flags \emph{is} an attack event, but the attached ATT\&CK id is usually wrong (Fig.~\ref{fig:confusion}). The
errors are \emph{systematic, not random}: the model confuses semantically adjacent siblings within a tactic ---
e.g.\ it labels all network-discovery (T1046) events as T1578, and maps account-manipulation (T1098) onto the
privilege-escalation cluster (T1548/T1548.005) --- yet it names \emph{distinctive} techniques reliably
(brute-force T1110 at $100\%$ exact). It identifies the malicious event and the right tactic family, then picks
the wrong specific sub-technique when several share observable signatures.

\begin{figure}[t]
\centering
\autofig{figures/confusion_A1.pdf}{Generate: \texttt{benchmark.plots --confusion <conf>.csv --arm A1}
(needs \texttt{benchmark.cli confusion --csv} from the run DB).}
\caption{Where technique attribution breaks (arm A1): per ground-truth technique, the
fraction the LLM labels exactly, with parent-only credit, with the wrong technique, or
misses entirely (real-AWS, qwen2.5:32b, arm A1).}
\label{fig:confusion}
\end{figure} Before the technique vocabulary was supplied, the model labelled events
with techniques outside the benchmark entirely (e.g., T1569, T1222); constraining the label space fixed the
out-of-vocabulary errors but not the within-vocabulary confusions.

\textbf{Rules (A4) failure.} The rules are exact on $16$ of $18$ techniques, failing only two sibling-ambiguity
cases --- \emph{Cloud Infrastructure Discovery} (T1580, mapped to the sibling T1526) and \emph{Resource
Hijacking} (T1496, mapped to T1578) --- which shows their high technique-level recall is a near-inverse of the
generator rather than general detection skill. They also \textbf{over-flag}: event-level precision falls to
$0.65$ on multi-stage (vs the LLM's $0.82$--$1.00$) and to zero on benign, because every
\texttt{ListBuckets}/\texttt{GetObject} trips a signature.

\textbf{Pre-filter recall stress.} In this run the pre-filter retained ground-truth events well --- $0.967$ on
multi-stage and $1.00$ on low-and-slow --- so the detections A3 recovers but A2 loses (Table~\ref{tab:twolens})
are dropped by the \emph{technique-signature} narrowing, not by wholesale removal of ground-truth events.

% =====================================================================
\section{Discussion}\label{sec:discussion}
% Owner: Anna
Our headline is a \emph{decomposition}: on AWS multi-stage kill chains, a mid-sized open LLM performs
cross-domain attack-event detection competitively and at higher precision than deterministic rules, but fails
at exact ATT\&CK-technique attribution. For a SOC, this suggests LLMs are promising as \emph{high-precision
triage / correlation} aids --- surfacing the events that matter with few false alarms --- while exact technique
labelling is better left to rules or human analysts. Multi-agent decomposition did not pay for its extra cost
(H2), and the pre-filter trades a modest cost saving for lost detections, so a cost-sensitive deployment might
skip it when the model can ingest the full stream.

\textbf{Why the rules win technique-level, and why that is not the whole story.} The rules are near-optimal
\emph{by construction}: their signatures are effectively the inverse of the scenario generator, so A4's
$0.87$/$0.97$ scores are an upper bound a real deployment would not reach. Notably, executing the scenarios
against real AWS did \emph{not} dissolve this advantage --- the captured telemetry is still the structured API
calls the benchmark issues, on which the signatures stay near-exact ($16$ of $18$ techniques). An independent
community Sigma rule-set (the SIGMA arm) lands between the extremes --- multi-stage technique-level recall
$0.74$ versus A4's $0.87$ and the LLM's $0.12$ --- confirming that A4's dominance reflects encoding the answer
key rather than a deliberately weak external baseline. The strict technique-level metric therefore structurally
favours rules that encode the answer key; we accordingly report event-level detection as the headline and
technique-level as a stricter secondary. A truly decisive test requires \emph{noisy} production telemetry ---
injected benign background traffic or a warmed managed-detector stream --- where no static rule can invert the
data; this is the primary follow-up. The choice of primary metric --- event- versus technique-level --- awaits
pre-registration sign-off.

% =====================================================================
\section{Threats to Validity}\label{sec:threats}
% Owner: shared (docs/week1/15)
\emph{External:} synthetic scenarios may not generalise---mitigated by real-incident
grounding and (partial) real-AWS validation. \emph{Authorship bias:} same team authors
scenarios and system---mitigated by a sealed held-out set and inter-rater review.
\emph{Construct:} bespoke metrics may reward conformity---mitigated by leading with
standard metrics. \emph{Internal:} LLM non-determinism---mitigated by a pinned model,
multi-seed runs, and reported variance. \emph{Statistical:} small per-category $n$---the
power analysis restricts confirmatory claims to two primary tests; other categories are
descriptive. \emph{Reproducibility:} model drift/deprecation---mitigated by version
pinning and full artifact release; adversarial LLM evasion is explicitly out of scope.
\emph{Baseline coverage.} We enabled Amazon GuardDuty on the sandbox account with a
$\sim$6-day baseline and collected its findings over the real-AWS run; it produced
\emph{no} findings mapping to any emulated technique (its only finding, root-credential
API usage, was unrelated account hygiene). We read this as a \emph{construct} limitation of
the telemetry rather than of GuardDuty: our scenarios execute the control-plane
\emph{shape} of attacks but not the signals GuardDuty is tuned for---the exfiltration path
is a synthetic-network fallback (no real VPC Flow Logs), the adversary IP addresses are
documentation-range rather than threat-intel-flagged, no crypto-mining instance actually
launches, and most stages are proxied, benign-looking API calls from a single legitimate
principal. We therefore exclude GuardDuty as a \emph{scored} baseline and retain the
community Sigma baseline (SIGMA arm) as the independent rules-only comparison, which
requires no managed service. Making a managed-detector comparison meaningful would require
raising scenario realism to emit GuardDuty-triggering signals (threat-intel-flagged
sources, real egress volume, genuinely launched instances)---a benchmark-realism upgrade we
leave to future work. LLMCloudHunter generates detection rules from CTI reports (a
different task) and is therefore positioned in related work rather than scored here.

% =====================================================================
\section{Conclusion}\label{sec:conclusion}
% Owner: Anna
We introduced CloudKC-Bench, the first open benchmark for AWS control-plane multi-stage
kill-chain reconstruction, and used it to run a pre-registered, four-arm controlled study
--- plus an independent community Sigma baseline --- on real-AWS telemetry. Our central
result is a \emph{decomposition} rather than a verdict: a mid-sized open LLM (qwen2.5:32b)
performs cross-domain attack-\emph{event} detection at higher precision than deterministic
rules (event-level precision up to $1.00$ vs $0.65$), yet fails at exact
ATT\&CK-\emph{technique} attribution (technique-level recall $\sim$0.12 vs the rules'
$0.87$), and its errors are \emph{systematic} confusions among sibling techniques rather
than random noise. We resolve both pre-registered hypotheses as negatives and report them
plainly: the LLM does not beat the rules on strict technique-level recall (H1), and
multi-agent decomposition adds no measurable value over a single agent (H2). The Sigma
baseline lands between the two, and an enabled GuardDuty produced no findings on our
telemetry --- both underscoring that the strict metric rewards signatures that encode the
answer key. For a SOC, the practical reading is that a mid-sized LLM is a promising
\emph{high-precision triage and correlation} aid, while exact technique labelling is better
left to rules or analysts. The durable contribution, however the metric debate settles, is
the benchmark itself: an open, ground-truthed, mechanically scored artifact that makes such
claims falsifiable.

\noindent\textbf{Future work.} Two threads follow directly from our results. First, closing
the \emph{signal gap} that made both the rules near-optimal and GuardDuty silent: injecting
benign background traffic and raising scenario realism so the telemetry carries the network-
and threat-intel-level signals that managed detectors and a noisy production stream would
surface --- the decisive test of whether the LLM's event-level precision advantage survives
outside a structured capture. Second, breadth: capable-model and larger real-AWS runs,
Zeek-derived network telemetry, the sealed held-out evaluation, and adversarial-evasion
scenarios.

\section*{Availability}
The generator, manifests, schema, matching function, and analysis scripts are released at
\url{https://github.com/therealannaa/CloudSentinel} \todo{swap to the archival DOI on release}.

\section*{Acknowledgment}
The authors thank Dr.~S.~Nagasundari (ISFCR, PES University) for supervision.

\bibliographystyle{IEEEtran}
\bibliography{references}

\end{document}
